%=================================宏包管理=====================================================
%	和配置有关的宏包在具体的配置区引用，这里只引用正文区用到的宏包
\usepackage{wallpaper}	% 封面背景包
\usepackage{amsmath,mathtools,amsthm,amsfonts,amssymb,bm}	% AMS包
%%%%%%%%%%%%%%%%%%%% 全局调整行间公式的上下间距 STRAT %%%%%%%%%%%%%%%%%%%%
% \AtBeginDocument{%
%   \setlength{\abovedisplayskip}{4pt plus 1pt minus 1pt}%
%   \setlength{\belowdisplayskip}{4pt plus 1pt minus 1pt}%
%   \setlength{\abovedisplayshortskip}{4pt plus 1pt}%
%   \setlength{\belowdisplayshortskip}{4pt plus 1pt minus 1pt}%
% }
\usepackage{etoolbox}
%\setlength{\abovedisplayskip}{3pt} %调节公式的上行之间的间距
 %   \setlength{\belowdisplayskip}{3pt} %调节公式的下行之间的间距
\newcommand{\setmydisplayskip}{%
  \setlength{\abovedisplayskip}{3pt plus 1pt minus 1pt}%
  \setlength{\belowdisplayskip}{3pt plus 1pt minus 1pt}%
  \setlength{\abovedisplayshortskip}{3pt plus 1pt minus 1pt}%
  \setlength{\belowdisplayshortskip}{3pt plus 1pt minus 1pt}%
}
% 在 \normalsize 之后追加我们的设置
\apptocmd{\normalsize}{\setmydisplayskip}{}{}
%%%%%%%%%%%%%%%%%%%% 全局调整行间公式的上下间距 END  %%%%%%%%%%%%%%%%%%%%

\usepackage{color}  %字体背景颜色包
\usepackage{wasysym}
% 使用：
%\CheckedBox      % ☑ 方框中带对号
%\Square          % ☐ 空方框
%\XBox            % ☒ 方框中带叉
%=================================页面边距=====================================================
%	geometry宏包使用教程：http://www.ctex.org/documents/packages/layout/geometry.htm
%	A4纸宽210mm，长297mm
%	left + right + textwidth = 210
%	top + bottom + textheight = 297
%	headheight：页眉文字高度，应当小于等于top
\usepackage{fancyhdr}  % 自定义页眉页脚
%\setlength{\headheight}{26pt}
%\addtolength{\topmargin}{-12pt}

\usepackage{geometry}    % 页面尺寸与边距
\usepackage{calc}        % 支持长度的算术运算
\usepackage{anyfontsize} % 任意字号支持（用于精确字号）
\geometry{%
	paperwidth  = 210mm,
	paperheight = 297mm,
	twoside,
	top         = 30mm,
	bottom      = 25mm,
	inner       = 28mm,
	outer       = 25mm,
	headheight  = 26pt,   % 页眉文字高（约5.1mm） 14.5
%	headsep     = 4.9mm,    % 10mm - 5.1mm ≈ 4.9mm（页眉顶距上缘20mm）
    headsep     =   8pt,    % 数值越小，页眉离正文越近。可以试 8pt、12pt、15pt 等找到合适的
	footskip    = 9.6mm,    % 使页脚顶部距下缘约17.5mm（需打印后校准）
}

%=================================页眉页脚=====================================================
%	fancy宏包使用教程：http://www.ctex.org/documents/packages/layout/fancyhdr.htm
%	fancypagestyle{样式名}可以自定义样式，并通过\pagestyle{样式名}和\thispagestyle{样式名}来使用
%   \leftmark可以获取不带星号的chapter标题内容，\rightmark可以获取到不带星号的section标题内容
%   L, C, R分别表示左中右，
%   E, O分别表示偶数页和奇数页
\usepackage{fontspec}
\setmainfont{Times New Roman} % 设置英文字体
\usepackage{setspace} % 段落行距包
\setstretch{1.5} % 设置行距为1.5倍行距
\raggedbottom

\usepackage{indentfirst} %段落首行缩进命令包
\setlength{\parindent}{2em}
\fancypagestyle{myfancy}{%
	\fancyhf{}	                                    % 清空所有定义
	\fancyfoot[CE,CO]{\thepage}                     % 设置页脚为当前页码
	\fancyhead[CE]{江苏科技大学本科毕业论文（设计）}        % 设置偶数页居中的页眉为江苏科技大学本科毕业论文（设计）
	\fancyhead[CO]{江苏科技大学本科毕业论文（设计）}
	%\fancyhead[CO]{\nouppercase{\leftmark}}         % 设置奇数页居中的页眉为当前章节名
	\renewcommand{\headrule}{%                      % 重定义headrule来实现双页眉装饰线效果
		\makebox[0pt][l]{\rule[.7\baselineskip]{\headwidth}{0.8pt}}%
	%	\rule[.6\baselineskip]{\headwidth}{0.4pt}\vskip-.8\baselineskip
	}
}

%=================================章节标题=====================================================
\ctexset{
	chapter = {
		name        = {},
		number      = \arabic{chapter},
		format      += \heiti\zihao{-3}\centering,  % 小三黑体居中
		beforeskip  = 15pt,   % 段前1.5行，设置30pt 
		afterskip   = 15pt,   % 段后1.5行
		fixskip     = true,
	},
	section = {
		format      += \heiti\zihao{4}\raggedright,  % 四号黑体居左
		beforeskip  = 10pt,   % 段前1.0行，设置20pt
		afterskip   = 10pt,
		fixskip     = true,
	},
	subsection = {
		format      += \heiti\zihao{-4}\raggedright, % 小四黑体居左
		beforeskip  = 10pt,   % 段前0.5行
		afterskip   = 10pt,
		fixskip     = true,
	},
	subsubsection = {
		format      += \heiti\zihao{-4}\raggedright, % 小四黑体居左(款)
		beforeskip  = 10pt,
		afterskip   = 10pt,
		fixskip     = true,
	},
}

%=================================目录设置=====================================================
%	titletoc宏包使用教程：https://blog.csdn.net/golden1314521/article/details/39926135
%						 https://blog.csdn.net/l_changyun/article/details/87431805
%
\usepackage{titletoc}
\usepackage{titlesec}
\renewcommand{\contentsname}{目\hspace{2em}录}
\setcounter{tocdepth}{2}      % 显示到 subsection(第3级)
\setcounter{secnumdepth}{3}
\ctexset{chapter/name = {第,章}}  % 章号显示"第X章"

%--------------------------------------------------------
% 一级(章 chapter)——顶格
%--------------------------------------------------------
\titlecontents%	章
{chapter}[4em]
{\heiti\zihao{-4}\vspace*{-4pt}}	% 章之间的上下间距  \bfseries
{\contentslabel{4em}}
{\hspace*{-4em}}
{~\titlerule*[0.6pc]{$.$}~\contentspage}


%%--------------------------------------------------------
%% 二级(节 section)——缩进一格
%%--------------------------------------------------------
\titlecontents%	节
{section}[3em]
{\zihao{-4}\vspace*{-4pt}}
{\contentslabel{2em}}
{\hspace*{-2em}}
{~\titlerule*[0.6pc]{$.$}~\contentspage}

%%--------------------------------------------------------
%% 三级(条 subsection)——缩进两格
%%--------------------------------------------------------
\titlecontents%	小节
{subsection}[6em]
{\zihao{-4}\vspace*{-4pt}}
{\contentslabel{3em}}
{\hspace*{-2em}}
{~\titlerule*[0.6pc]{$.$}~\contentspage}
 

%=================================参考文献=====================================================
%   biblatex宏包使用教程：
%	https://www.overleaf.com/learn/latex/Bibliography_management_with_biblatex
%
\usepackage[%
backend=biber,				% 设置使用biber进行编译，也可以使用bibtex，但是功能更少
style=gb7714-2015,			% 设置风格样式为国家标准gb7714-2015
sorting=none,%设置排序按照年份，名字，标题进行排序，若想按照引用顺序排序，将其设置为none即可
gbnamefmt=lowercase           % 将作者名全部大写改为仅首字母大写
]{biblatex}       			% 参考文献包
\addbibresource{ref.bib}    % 加载参考文献的文件

%%参考文献工具，加载biblatex宏包，,其后端backend使用biber
%\usepackage[backend=biber,style=gb7714-2025,gbfootbib=true,gbalign=gb7714-2025]{biblatex}
%biblatex宏包的参考文献数据源即bib文件加载方式
%\addbibresource[location=local]{references.bib}

%=================================定理环境=====================================================
% 	自定义定理类环境（定义，引理，定理，推论，例，注）
% 		定理环境命令：\newtheorem{name}[counter]{text}[section]
% 			name:		标识这个环境的关键字（用于编程）
%			counter:	（可选）编号计数器，默认使用自己的计数器，可以传入其他环境的name来共享计数器
% 			text:		真正在文档中打印出来的定理环境的名字
% 			section:	（可选）定理编号依赖的某个章节层次，默认不依赖。
%
\newtheorem{theorem}{\hskip 2em{定理}}[chapter]
\newtheorem{definition}[theorem]{\hskip 2em{定义}}
\newtheorem{lemma}[theorem]{\hskip 2em{引理}}
\newtheorem{corollary}[theorem]{\hskip 2em{推论}}
% 例，注各自独立编号，无需考虑编号共享的问题，直接创建。证明关键词加粗
\newtheorem{example}{\hskip 2em{例}}[chapter]
\newtheorem{remark}{\hskip 2em{注}}[chapter]
\renewcommand{\proofname}{\hskip 2em \bf 证明}

% 去掉定理后面的小点，不建议使用（默认注释掉）
% \usepackage{xpatch}
% \makeatletter
% \AtBeginDocument{\xpatchcmd{\@thm}{\thm@headpunct{.}}{\thm@headpunct{}}{}{}}
% \makeatother

% 公式按section编号，若想按chapter编号，注释掉这条即可
% \numberwithin{equation}{section}

%=================================图表环境=====================================================
% enumitem宏包设置参考自江苏科技大学学位论文模板
\usepackage[inline]{enumitem}					% 列表工具包
\usepackage{graphicx,ragged2e}					% 插图工具包
\usepackage{subcaption}							% 子图标题包
\usepackage{bicaption}							% 图片标题包
\setlist{%	设置列表样式
	topsep=0.3em, 			% 列表顶端的垂直空白
	partopsep=0pt, 			% 列表环境前面紧接着一个空白行时其顶端的额外垂直空白
	itemsep=0ex plus 0.1ex, % 列表项之间的额外垂直空白
	parsep=0pt, 			% 列表项内的段落之间的垂直空白
	leftmargin=1.5em, 		% 环境的左边界和列表之间的水平距离
	rightmargin=0em, 		% 环境的右边界和列表之间的水平距离
	labelsep=0.5em, 		% 包含标签的盒子与列表项的第一行文本之间的间隔
	labelwidth=2em 			% 包含标签的盒子的正常宽度；若实际宽度更宽，则使用实际宽度。
}
\graphicspath{figure/}		% 设置图片存放目录
\usepackage{longtable,booktabs} % 跨页长表格
 

%=================================代码环境=====================================================
% 使用listings宏包来插入代码
\usepackage{listings}	% 代码环境包
\renewcommand{\lstlistingname}{代码}	% 重命名代码块标题为算法，例如：算法1.2
%\renewcommand{\lstlistingname}{算法}	% 重命名代码块标题为算法，例如：算法1.2
%\lstset{% 设置算法样式
%	keywordstyle=\bfseries, % 设置关键词加粗
%	basicstyle=\ttfamily, 	% 设置基础样式字体为等宽
%	commentstyle=\ttfamily, % 基本和注释的字体都使用默认的等宽，而非texlive调用的中文字体
%	showstringspaces=false, % 不显示中间的空格
%	breaklines=true,  		% 对过长的代码自动换行
%	frame=single  			% 边框
%    basicstyle=\small\ttfamily, % 基本字体样式（小号等宽字体）
%    keywordstyle=\color{blue}\bfseries, % 关键字样式（蓝色加粗）
%    commentstyle=\color{green!60!black}\itshape, % 注释样式（深绿斜体）
%    stringstyle=\color{red},    % 字符串样式（红色）
%    showstringspaces=false,     % 不显示字符串中的空格
%    breaklines=true,            % 自动换行
%    frame=single,               % 添加边框
%    tabsize=4                   % Tab 宽度为 4 个字符
%}


% Algorithms 
\usepackage[plainruled,linesnumbered,algochapter]{algorithm2e}
%=================================引用工具=====================================================
% hyperref宏包教程https://www.jianshu.com/p/58e7d0a6d97a
% 实现超链接功能
\usepackage{hyperref}	% 交叉引用包
\hypersetup{%	设置交叉引用属性
	colorlinks=true,	% 设置可跳转的链接为颜色，而不是方框
	urlcolor=blue,		% 设置各种链接的颜色均为黑色
	linkcolor=black,
	anchorcolor=black,
	citecolor=black
}

\makeatletter
% ===== 基本中英文名称 =====
\newcommand{\algoenname}{\textbf{Algo.}}        % 算法英文标号前缀
\renewcommand{\algorithmcfname}{算法}            % 算法中文名称

% ===== 间距设置 =====
\setlength\AlCapSkip{1.2ex}
\SetAlgoSkip{1pt}
\setlength{\algoheightruledefault}{1.5pt}

% ===== 标题计数器格式（带章节号）=====
\expandafter\ifx\csname algocf@within\endcsname\relax
  \renewcommand\thealgocf{\@arabic\c@algocf}
\else
  \renewcommand\thealgocf{\csname the\algocf@within\endcsname-\@arabic\c@algocf}
\fi

% ===== 标题整体居中（中英文各占一行）=====
\renewcommand{\algocf@makecaption}[2]{%
  \addtolength{\hsize}{\algomargin}%
  \hskip .5\algomargin%
  \parbox[t]{\hsize}{%
    \centering
    \setlength{\baselineskip}{1.5em}%
    \AlCapFnt{}#1\algocf@typo~\AlCapFnt{}#2%
  }%
  \addtolength{\hsize}{-\algomargin}%
}

% ===== 双语标题命令 =====
% #1 中文标题  #2 英文标题
\newcommand{\AlgoBiCaption}[2]{%
  \caption[#1]{#1\\ \algoenname\ \thealgocf~~#2}%
  \addcontentsline{aen}{algoen}{\protect\numberline{ }{#2}}%
}

% ===== ruled 线宽 =====
\setlength{\algoheightruledefault}{1.5pt}
\def\@algocf@post@ruled{\kern\interspacealgoruled\hrule height1.5pt\relax}%
\makeatother

%======== 中文关键字 ========%
\SetKwInput{KwIn}{输入}
\SetKwInput{KwOut}{输出}


%=========== Julia代码设置 ============================
% 使用Julia的格式输入，需要导入
\usepackage{jlcode}
% 使用Julia的格式输入，或者是另外一个头文件,
% \input{settings/julia.jl}
% \begin{jllisting}
% \end{jllisting}

%=================================其余设定=====================================================
% 重新定义一些常用的数学符号
\renewcommand{\Re}{\operatorname{Re}}
\renewcommand{\Im}{\operatorname{Im}}
\newcommand{\mi}{\mathrm{i}}
\newcommand{\md}{\mathrm{d}}
\newcommand{\me}{\mathrm{e}}
 \DeclareMathOperator{\argmin}{argmin}
\DeclareMathOperator{\supp}{supp}
\DeclareMathOperator{\prox}{prox}
\DeclareMathOperator{\sign}{sign}


